\section{Introduction}
\label{sec:introduction}

The 2025 NeurIPS / ICML cycle established a dominant template for reinforcement-learning-driven biological sequence design: take a pretrained generative model (a protein language model, a discrete diffusion model, or a DNA-LLM); fine-tune with PPO/DPO/GRPO/KTO against a fitness proxy; add a naturalness regulariser. Methods include $\delta$-Conservative Search \citep{kim2025deltacs}, $\mu$Protein \citep{sun2025muprotein}, the steering work of \citet{yang2025steering}, GLID$^2$E \citep{cao2025glid2e}, BioReason \citep{bioreason2025}, and the Bayesian active-learning toolkit ALDE \citep{yang2025alde}. They all share a single architectural choice: an \emph{undifferentiated} policy gradient against a scalar reward.

This pooling is the central pathology of RL-for-biology in 2025. When the proxy reward becomes unreliable on out-of-distribution sequences, the policy hacks the proxy. Existing mitigations (e.g.\ $\delta$-CS) clip the exploration radius using proxy variance, treating the symptom (mis-calibrated rewards) rather than the cause (a gradient estimator that cannot tell us whether the policy is exploiting known regions or shifting to new ones).

The cause is structural. The policy gradient
\(
\nabla_\theta J(\theta) = \mathbb{E}_{x \sim \pi_\theta}\!\left[ R(x)\,\nabla_\theta \log \pi_\theta(x) \right]
\)
pools two effects: it up-weights trajectories the policy already supports (the analogue of evolutionary \emph{selection}), and it shifts the policy's support to new regions (the analogue of \emph{transmission}, i.e.\ mutation/inheritance). \citet{price1970selection} decomposed any change in mean trait value across an evolving population into precisely these two terms, exactly. \citet{frank2025price} recently proved that the algebraic form $\Delta\theta = M f + b + \xi$ --- the force--metric--bias law --- subsumes natural selection, Bayesian updating, SGD, and policy gradients as algebraically identical objects, with the Price decomposition lifting to all of them. Crucially, that paper made no algorithmic claim. We do.

\paragraph{Contributions.} (i) A new RL primitive: estimate the selection and transmission components of $\nabla J$ separately on policy, partitioned by membership in the previous policy's support. (ii) An adaptive controller that drives the empirical Price ratio $\rho_t = \|g_S\|/(\|g_S\|+\|g_T\|)$ towards a target $\rho^*_t$ derived from the reward landscape's autocorrelation length. (iii) A new diagnostic: $\rho_t$ itself, used as an early-warning signal for reward hacking. (iv) Theorem~1 (exactness for factorised categorical policies), empirically validated to cosine $1.0000$ on every round of every NK seed. (v) Empirical results on 149{,}361-variant GB1 (Wu 2016), the 200-label NK reward-hacking experiment, and the seven-K NK theory sweep. (vi) Public release of the code under MIT licence.

\paragraph{Where this is and is not novel.} It is \emph{not} novel to use RL for biological sequence design \citep{angermueller2020dyna,jain2022biological,lee2024latprotrl,kirjner2024smoothed}; it is \emph{not} novel to be ruggedness-aware \citep{towers2024ruggedness,sun2025muprotein}; it is \emph{not} novel to relate evolution and learning at a verbal level \citep{salimans2017evolution}. What is novel is the formal operationalisation of \citeauthor{frank2025price}'s 2025 unification as an algorithmic primitive, the Price-ratio diagnostic, and the regret-by-autocorrelation framing. The paper does not propose a better protein optimiser; it proposes that the algebraic structure of the optimiser was wrong, and shows that biology is the domain where the rewriting pays off most cleanly.
